\chapter{Durchführung} 
\label{ch:Durchfuehrung}

\section{Messverfahren}
\label{sec:Messverfahren}

Der Versuch gliedert sich in drei Hauptteile, welche jeweils unterschiedliche Aspekte der elektromagnetischen Induktion beleuchten. Im ersten Versuchsteil wurde die Helmholtzspule an eine Gleichspannungsquelle angeschlossen, die gleichzeitig den Strom misst. Die Induktionsspule wurde mittels eines Motors in Rotation versetzt, wobei die Frequenz des Motors über die Betriebsspannung variiert wurde. Die Frequenz der Induktionsspannung wurde am Oszilloskop direkt abgelesen. Mit diesem Setup wurden die Messungen zum Induktionsgesetz aufgenommen, indem sowohl die Frequenzabhängigkeit bei konstanter Stromstärke als auch die Stromstärkeabhängigkeit bei konstanter Frequenz untersucht wurden.

Im zweiten Versuchsteil wurde mit einem Funktionengenerator eine Wechselspannung an die Helmholtzspulen angelegt. Mittels Oszilloskop wurden die Peak-to-Peak-Spannung und die Frequenz gemessen. Zusätzlich wurde ein Multimeter in Reihe geschaltet, um die effektive Stromstärke zu messen. Die Induktionsspannung der nun ruhenden Spule wurde weiterhin am Oszilloskop abgelesen. Hierbei wurde systematisch der Winkel zwischen der Normalen der Spulenfläche und dem Magnetfeld verändert, um die Winkelabhängigkeit zu untersuchen. Weiterhin wurde das Transformationsverhalten des Aufbaus als Lufttransformator analysiert, indem das Verhältnis zwischen Induktionsspannung und Primärspannung über ein breites Frequenzspektrum vermessen wurde. Aus dem Frequenzgang des Widerstands der Helmholtzspule konnte zudem ihre Induktivität bestimmt werden.

Im dritten und letzten Versuchsteil wurde die Induktionsspule erneut in Drehung versetzt. Der Versuchsaufbau wurde soweit möglich nach Norden ausgerichtet, allerdings erfolgte die Bestimmung der Nordrichtung hauptsächlich anhand der Orientierung des Gebäudes, da kein zuverlässiger Kompass vorhanden war. In dieser Ausrichtung wurde die Peak-to-Peak-Induktionsspannung am Oszilloskop ohne externe Kompensation gemessen. Anschließend wurde die Helmholtzspule mit Gleichstrom angeschaltet und deren Spannung so variiert, bis die Induktionsspannung minimal wurde. Diese Kompensationsmethode erlaubt die Trennung von horizontaler und vertikaler Komponente des Erdmagnetfeldes. Aufgrund des verrauschten Signals wurde ein Tiefpassfilter zwischengeschaltet, welcher zwar das Signal glättete, aber zugleich die gemessene Induktionsspannung signifikant reduzierte. Die Bestimmung war insgesamt ungenau, da die Anzeigewerte der Induktionsspannung stark schwankten.

\section{Versuchsaufbau}
\label{sec:Versucshaufbau}

Der verwendete Versuchsaufbau besteht im Kern aus einer Helmholtzspulenanordnung mit einer im Zentrum drehbar gelagerten Induktionsspule. Die Helmholtzspule verfügt über einen Durchmesser von 295 mm, einen Spulenabstand von 147 mm und jeweils 127 Windungen pro Spule. Die Induktionsspule weist 4000 Windungen auf und hat eine Querschnittsfläche von 41.7 cm$^2$. Zur Spannungs- und Strommessung kommen ein Oszilloskop, ein Multimeter sowie diverse Netzteile zum Einsatz. Ein Leistungsfunktionsgenerator ermöglicht die Einspeisung von Wechselspannungen variabler Frequenz, während ein Antriebsmotor mit Treibriemen die rotation der Induktionsspule übernimmt.

Zur Filterung des verrauschten Signals im Abschnitt zur Erdmagnetfeldmessung wurde ein RC-Tiefpassfilter in den Messkreis integriert. Alle Geräte wurden so verschaltet, dass sowohl Gleichspannungsquellen für statische Magnetfelder als auch Wechselspannungsquellen für periodische Feldänderungen genutzt werden konnten. Die gesamte Anordnung befindet sich innerhalb einer Schutzumrandung, welche eine sichere Bedienung gewährleistet. Eine grobe Ausrichtung nach geografischem Norden erfolgt visuell entlang bekannter Gebäudeachsen, wobei aufgrund fehlendem Kompass hier eine systematische Unsicherheit bestehen bleibt.